\documentclass[aps,pra,reprint,groupedaddress,floatfix]{revtex4-2}
\usepackage{graphicx}
\usepackage{amsmath,amssymb}
\usepackage{placeins}
% Float tuning (reprint two-column carries two large figure*s: D04 2x2, D06 trio).
% Legibility-first: keep panel sizes, relax placement fractions instead.
\renewcommand{\topfraction}{0.9}
\renewcommand{\bottomfraction}{0.8}
\renewcommand{\textfraction}{0.07}
\renewcommand{\floatpagefraction}{0.6}

\graphicspath{{figures/}}

\begin{document}

\title{Model and finite-size properties of the SSH XXZ chain\\
with a normalized topological reflection marker}

\author{Xue Qiu}\thanks{These authors contributed equally.}
\author{Yijie Han}\thanks{These authors contributed equally.}
\author{Jinlong Yu}
\affiliation{School of Physics and Optoelectronic Engineering, Hainan University, Haikou 570228, China}

\date{\today}

\begin{abstract}
Variational eigensolvers have become a cornerstone of near-term quantum
simulation, yet topological phases remain hard to certify on hardware
because theorists' preferred markers are rarely directly measurable.
In this work, we prepare the competing trivial, topological, and
antiferromagnetic phases of the dimerized XXZ chain with
symmetry-preserving orbit ansatzes built on phase-matched initial
states. The finite-size phase diagram at $L=8$ is
mapped with the normalized reflection marker $\tilde{Z}_{\mathcal R}$,
resolving the three phases, with gap scaling
confirming controlled finite-size effects. State preparation uses only
directly measurable markers---the AFM structure factor $S(\pi)$ and the
string operator---while premium-qubit screening and tensor-product
readout mitigation bring the protocol to superconducting hardware.
Hardware and simulator trends
are compared on equal footing via internal $0$--$1$ normalization,
across two dimerizations and both directly measurable markers.
\end{abstract}

\maketitle

\section{Introduction}
\label{sec:intro}

Variational quantum eigensolvers have become a cornerstone of
near-term quantum simulation, trading circuit depth for a classical
optimization loop. For topological phases, however, the standard
variational workflow faces an extra hurdle: the order parameters that
theorists prefer are often the hardest to measure on hardware.

Topological phases on real hardware have been pursued both by direct
construction and by variational optimization. On superconducting
processors, Satzinger~\textit{et al.}~\cite{satzinger2021realizing}
prepared toric-code ground states and measured topological entanglement
entropy, demonstrating topological order itself as the target; direct
constructions of this kind yield satisfactory fidelities, however, they
are only applicable to exactly solvable fixed-point states.
For variational approaches, Sun~\textit{et al.}~\cite{sun2023efficient}
designed initialization-plus-variational ansatzes for symmetry-protected
topological phases and verified string order and edge modes on IBM
devices, although the ansatz structure has largely been determined by
trial and error for each target phase.

While the main idea of phase-matched initialization is rather simple,
its full potential for generic non-integrable models---with a single
ansatz family covering competing orders and with markers that survive
the transition to hardware---has not been uncovered.
The neighboring literatures we draw on are
established: one-dimensional symmetry-protected phases are classified by
projective representations and matrix-product
states~\cite{chen2011classification,schuch2011classifying};
randomized measurements reconstruct entanglement and many-body
invariants from few
settings~\cite{huang2020predicting,elben2022toolbox};
variational algorithms and their trainability limits are
surveyed~\cite{cerezo2021variational,tilly2022variational,mcclean2018barren};
error mitigation spans probabilistic cancellation and zero-noise
extrapolation~\cite{temme2017error,kandala2019error}; and variational
hardware experiments range from molecular spectra to spin
chains~\cite{colless2018robust,yu2023simulating}.
The two diagnostics we use have dedicated literatures of their own.
The dynamical spin structure factor of the anisotropic spin-$1/2$
Heisenberg chain was computed by Pereira~\textit{et
al.}~\cite{pereira2006dynamical}, establishing the momentum-resolved
response whose static limit underlies our $S(\pi)$ AFM marker.
String order was introduced for valence-bond
phases~\cite{dennijs1989preroughening}, characterized symmetry-wise for
quantum spin lattices~\cite{perezgarcia2008string}, and observed via
correlated particle-hole pairs in low-dimensional Mott
insulators~\cite{endres2011observation}; our string operator is the
directly measurable counterpart tracked across the same sweep.
On hardware, limited connectivity makes qubit mapping a first-class
problem~\cite{li2019tackling}; our stabilizer-plus-readout screening
answers the same question measurement-first, ranking the subgraphs the
experiment will actually use. Readout mitigation spans correlated
unfolding~\cite{bravyi2021mitigating,nation2021scalable} and model-free
expectation-value correction~\cite{vandenberg2022modelfree}; our
per-qubit $2\times2$ unfolding with nonnegativity sits in this
tensor-product family by construction. Finite-size gap scaling is
standard practice in spin-chain
numerics~\cite{sandvik2010computational}, with the Haldane-gap problem
as its classic case~\cite{haldane1983nonlinear,affleck1987rigorous,affleck1989quantum};
our $L=8$ assignments with gap-scaling checks follow that discipline,
with full curves in the Appendix.
In this work, we extend phase-matched variational preparation to the
dimerized XXZ chain, where trivial, topological, and antiferromagnetic
phases compete, and we characterize the prepared states with directly
measurable markers ($S(\pi)$, string) alongside the
reflection invariant $\tilde{Z}_{\mathcal R}$.
Our approach extends phase-matched preparation to a setting where three
orders compete within one model, while accounting for hardware
constraints from the start: premium-qubit screening, readout mitigation,
and markers restricted to directly measurable correlators.
Whether the prepared states track the simulator trends on hardware is
decided by the comparison in Sec.~\ref{sec:comparison}.

The paper is organized as follows. We first introduce the model and
its symmetries (Sec.~\ref{sec:model}) and map the finite-size phase
diagram (Sec.~\ref{sec:phasediagram}), with gap scaling in the Appendix.
We then define the directly measurable phase markers
(Sec.~\ref{sec:markers}) and prepare the variational states
(Sec.~\ref{sec:ansatz}), bring the protocol to hardware
(Sec.~\ref{sec:hardware}), and compare hardware with simulator results
(Sec.~\ref{sec:comparison}).

\FloatBarrier
\section{Model and symmetries}
\label{sec:model}

We consider the dimerized XXZ chain with $L$ sites and open boundary
conditions, sites labeled $1,\dots,L$. The dimerized alternating-bond
structure goes back to Su, Schrieffer, and Heeger~\cite{su1979solitons}. With $X_m,Y_m,Z_m$ the Pauli
operators on site $m$, the Hamiltonian reads
\begin{equation}
H(s,\delta)=H_o+H_e,
\end{equation}
where $H_o$ ($H_e$) couples odd (even) bonds, $s$ is the dimerization
parameter, and $\delta$ controls the anisotropy:
\begin{equation}
H_o=(1-s)\sum_{j=1}^{L/2}h_{2j-1,2j},\qquad
H_e=s\sum_{j=1}^{L/2-1}h_{2j,2j+1},
\end{equation}
\begin{equation}
h_{kl}=e^{-\delta}(X_kX_l+Y_kY_l)+e^{+\delta}Z_kZ_l,
\end{equation}
where $h_{kl}$ is the two-site bond operator with
$J_{xx}=e^{-\delta}$ and $J_z=e^{+\delta}$; $\delta=0$ is the Heisenberg
isotropic point, while $\delta>0$ leans toward Ising.
Unless stated otherwise we take $L=8$.

$H$ commutes with the total magnetization
$Z_{\rm tot}=\sum_m Z_m$, the global spin flip $\prod_m X_m$, and the
full-chain reflection $R$ ($m\leftrightarrow L+1-m$), so evolution under
$H$ never leaves the initial symmetry sector. For $\delta\ne0$ the symmetry is $U(1)$ ($Z_{\rm tot}$
conserved); at $\delta=0$ it enlarges to $SU(2)$.

The topological phase here belongs to the one-dimensional
symmetry-protected family whose study began with Haldane's conjecture
of a finite gap for integer-spin antiferromagnetic
chains~\cite{haldane1983nonlinear,dennijs1989preroughening} and the
exact valence-bond-solid ground state found at a special point of the
spin-$1$ chain~\cite{affleck1987rigorous,dennijs1989preroughening}.
That phenomenology is now classified: gapped symmetric phases in one
dimension are labeled by the second cohomology group, i.e.\ projective
representations of the symmetry~\cite{chen2011classification}, equivalently
in matrix-product-state
form~\cite{schuch2011classifying}, with the degeneracy of the
entanglement spectrum as the operational
fingerprint~\cite{pollmann2010entanglement,li2008entanglement}.
Our dimerized XXZ chain realizes three competing orders---trivial,
topological, and antiferromagnetic---within a single tunable model,
which is why a marker that resolves all three at once is needed.

To distinguish the phases we use the normalized topological reflection
marker, following the partial-reflection many-body invariant of
Elben~\textit{et al.}~\cite{elben2020manybody}. Let $\rho_I$ be the reduced density matrix of the ground state
on the central block $I$ (for $L=8$, sites $3$--$6$) and
${\mathcal R}_I$ the mirror exchange about the central bond. With
$Z_{\mathcal R}={\rm Tr}(\rho_I{\mathcal R}_I)$ and the corresponding
quantities $\rho_{I_{1,2}}$ on the left/right halves of $I$,
\begin{equation}
\tilde{Z}_{\mathcal R}=
\frac{Z_{\mathcal R}}
{\sqrt{[{\rm Tr}(\rho_{I_1}^2)+{\rm Tr}(\rho_{I_2}^2)]/2}},
\end{equation}
where the denominator is the mean purity of the two halves.
In other words, $\tilde{Z}_{\mathcal R}$ is the reflection expectation
normalized by the local mixedness, so that it approaches $+1$ ($-1$) in
the trivial (topological) phase and $0$ in the magnetically ordered
phase, independent of the overall entanglement level.
The finite-size phase diagram mapped with this marker is reported in
Sec.~\ref{sec:phasediagram}.

\FloatBarrier
\section{Finite-size phase diagram}
\label{sec:phasediagram}

The normalized marker over the $99\times99$ grid $s=i/100$,
$\delta=j/50$ ($i,j=1,\dots,99$) is reported in Fig.~\ref{fig:D01}.
Each grid point is computed independently from the stored ground state,
without interpolation or smoothing; degenerate points keep the solver's
lowest eigenvector and carry a degeneracy flag in the dataset.
The resulting map defines the parameter-space regions used to select
representative points and to specify the variational start/end protocol.
The marker is used here as a phase-labeling convention; the finite-size
drift of the boundaries is analyzed separately in the Appendix.

\begin{figure}[htbp]
\centering
\includegraphics[width=\columnwidth]{exp01_D01.pdf}
\caption{Normalized marker $\tilde{Z}_{\mathcal R}(s,\delta)$ at $L=8$ on the
$99\times99$ grid, computed point by point without interpolation.
\label{fig:D01}}
\end{figure}

\FloatBarrier
\section{Phase markers}
\label{sec:markers}

\subsection{Structure factor and the AFM label}
\label{sec:sq}

To identify the antiferromagnetic phase we use the structure factor
\begin{equation}
S(q)=\frac{1}{L}\sum_{i,j}e^{iq(i-j)}\langle Z_iZ_j\rangle,
\end{equation}
where $\langle Z_iZ_j\rangle$ is the ground-state correlator and $L=8$
throughout this section.
That is, $S(q)$ is the Fourier weight of N\'eel-type correlations, so
long-range antiferromagnetic order appears as a Bragg-like peak at
$q=\pi$. The $S(q)$ curves at one representative point in each
of the three phases (labels trivial / topo / AFM, coordinates in the
dataset) on the shared $q$ grid ($99$ points including $q=\pi$) are shown in
Fig.~\ref{fig:D03}.
The $q=\pi$ component is used as the AFM diagnostic throughout the
hardware-facing protocol. Its empirical phase discrimination is reported
with the results rather than assumed as part of the measurement method.

\begin{figure}[htbp]
\centering
\includegraphics[width=\columnwidth]{exp02_D03.pdf}
\caption{Structure factor $S(q)$ at three phase-representative points (one
per phase) on the shared $q$ grid; only the AFM point peaks at $q=\pi$.
\label{fig:D03}}
\end{figure}

\FloatBarrier
\subsection{Topological heatmaps}
\label{sec:heatmaps}

Four heatmaps over the same $(s,\delta)$
grid are collected in Fig.~\ref{fig:D04}: (a) the AFM structure factor $S(\pi)$; (b) the string operator,
which distinguishes symmetry-protected topological order since the
work of den Nijs and Rommelse~\cite{dennijs1989preroughening};
(c) the ZR-like combination $Q=\frac{4}{3}+2O_{\rm str}-\frac{1}{6}S(\pi)$,
where $O_{\rm str}$ is the string expectation; and
(d) the normalized marker $\tilde{Z}_{\mathcal R}$ of
Ref.~\cite{elben2020manybody}, recomputed here independently of Fig.~\ref{fig:D01}.
The four quantities are evaluated on the same parameter grid so that
their diagnostic roles can be compared consistently in the results.
The string signal in panel (b) is the long-range limit of the correlator
\begin{equation}
O^z_{\rm str}(i,j)=\Big\langle S^z_i\exp\!\Big(i\pi\!\sum_{k=i+1}^{j-1}S^z_k\Big)S^z_j\Big\rangle,
\end{equation}
where $S^z_k$ is the spin-$z$ operator at site $k$: the phase string
detects hidden antiferromagnetic order that no local two-point function
sees, the order parameter introduced for valence-bond
phases~\cite{dennijs1989preroughening} and now understood through the
symmetry-protected degeneracy of the entanglement
spectrum~\cite{pollmann2010entanglement}.
Entanglement-based fingerprints of this kind began with the
entanglement spectrum as a topological diagnostic~\cite{li2008entanglement}
and the universal constant term $S=\alpha L-\gamma$ of the entanglement
entropy~\cite{kitaev2006topological,levin2006detecting}; our panels (a)
and (b) trade that full entanglement resolution for correlators
measurable with single-qubit readouts.
Indeed, randomized measurements now reconstruct such quantities from
few settings: classical shadows predict $M$ linear properties from
$N\ge{\cal O}(\log M\,\max\|O_i\|^2_{\rm shadow}/\epsilon^2)$
shots~\cite{huang2020predicting}, derandomized into deterministic Pauli
measurements~\cite{huang2021derandomization}, biased toward the target
Hamiltonian~\cite{hadfield2022measurements}, extended to multi-shot
estimators~\cite{zhou2023performance}, mutually unbiased
bases~\cite{wang2023classical}, and even channel
tomography~\cite{kunjummen2023shadow}, as
surveyed in Ref.~\cite{elben2022toolbox}; second R\'enyi entropies and
quantum Fisher information have been measured this
way~\cite{brydges2019probing,vitale2024robust}.
In contrast to $\tilde{Z}_{\mathcal R}$, panels (a) and (b) involve only
standard Pauli correlators accessible to direct measurement.
Since $\tilde{Z}_{\mathcal R}$ is not directly measurable on real
hardware, the hardware-facing protocol uses state preparation together
with the directly measurable quantities in panels (a) and (b).

\begin{figure*}[p]
\centering
\begin{minipage}[b]{0.48\textwidth}
\centering
\includegraphics[width=\textwidth]{exp02_D04a.pdf}\\(a)
\end{minipage}
\hfill
\begin{minipage}[b]{0.48\textwidth}
\centering
\includegraphics[width=\textwidth]{exp02_D04b.pdf}\\(b)
\end{minipage}
\\
\begin{minipage}[b]{0.48\textwidth}
\centering
\includegraphics[width=\textwidth]{exp02_D04c.pdf}\\(c)
\end{minipage}
\hfill
\begin{minipage}[b]{0.48\textwidth}
\centering
\includegraphics[width=\textwidth]{exp02_D04d.pdf}\\(d)
\end{minipage}
\caption{Topological heatmaps over the $(s,\delta)$ grid of Fig.~\ref{fig:D01}:
(a) $S(\pi)$, (b) string operator,
(c) $Q$, (d) $\tilde{Z}_{\mathcal R}$. \label{fig:D04}}
\end{figure*}

\FloatBarrier
\section{Variational state preparation}
\label{sec:ansatz}

\subsection{Illustrative start and end points}
\label{sec:points}

For each phase we define an illustrative start/end pair on the
$\tilde{Z}_{\mathcal R}$ map (Fig.~\ref{fig:D05}). The end points are
classical reference targets chosen deep inside the corresponding phase;
they are demonstration endpoints for the variational protocol, not the
actual final evolution endpoints. Keeping the reference points away from
the phase boundaries reduces sensitivity to finite-size boundary drift.
Start points share the dot style in three colors; end markers differ in
both shape and color, with arrows showing the reference direction of
each pair.

\begin{figure}[htbp]
\centering
\includegraphics[width=\columnwidth]{exp01_D05.pdf}
\caption{Illustrative start/end pairs on the $\tilde{Z}_{\mathcal R}$
phase diagram of Fig.~\ref{fig:D01}, one pair per phase with the
reference end point deep inside it; colors distinguish pairs and arrows
mark reference direction.
\label{fig:D05}}
\end{figure}

\subsection{Variational circuit structure}

Our variational circuits follow the variational eigenvalue solver
paradigm~\cite{peruzzo2014variational}: prepare an initial state, then
apply a parameterized ansatz. The classical loop minimizes
\begin{equation}
E(\theta)=\langle\psi(\theta)|H|\psi(\theta)\rangle,
\end{equation}
where $|\psi(\theta)\rangle$ is the circuit output at parameters
$\theta$. Variational algorithms of this form are surveyed in
Refs.~\cite{cerezo2021variational,tilly2022variational}, which collect
cost-function design, ansatz families, and best practices.
In practice the optimization landscape, not the
expressivity, is the bottleneck: gradients vanish exponentially with
qubit number in random circuits forming
$2$-designs~\cite{mcclean2018barren}, hardware-efficient circuits exhibit
abrupt trainability transitions with depth~\cite{campos2021abrupt}, while
optimal parameters can concentrate inverse-polynomially in problem
size~\cite{akshay2021parameter}, smooth Hamiltonian-variational solutions
transfer across sizes with non-vanishing
gradients~\cite{mele2022avoiding}, and barren-plateau regimes can hide
exponentially many traps~\cite{nemkov2025barren}. Krylov-like iterative
solvers offer a complementary route to ground
states~\cite{bharti2021iterative}. Hardware-efficient variational
optimization was first demonstrated on superconducting hardware for
small molecules and quantum
magnets~\cite{kandala2017hardware}, and spin-chain variational simulation
has since been pushed to $102$ qubits on cloud superconducting
devices~\cite{yu2023simulating}. Our orbit ansatz differs from all of
these in tying parameters by reflection symmetry and fixing the
entanglement structure at initialization, as described next.
The three initial states (odd-bond singlet
product and its counterparts, all in the $P=-2$ sector, where $P$ labels
the combined symmetry eigenvalue) pair with an
orbit ansatz in which mirror-symmetric bond pairs share rotation angles
(odd bonds $O_j\leftrightarrow O_{M+1-j}$, even bonds
$E_j\leftrightarrow E_{M-j}$). That is, the ansatz is built in two steps:
first fix the entanglement structure with a phase-matched initial state,
then dress all bonds with rotations tied by the reflection symmetry, so
the parameter count grows with the number of bond orbits rather than
bonds. One schematic per phase is shown in Fig.~\ref{fig:D06}(a)--(c): initial state
plus required ansatz circuit for the trivial, topological, and
AFM cases.
The construction keeps the trial states in the target symmetry sector.
The quality of the resulting states is evaluated later by the
hardware-simulator benchmarks.

\begin{figure}[htbp]
\centering
\includegraphics[width=\columnwidth]{exp04_D06a.pdf}
\caption{(a) Trivial initial state plus required ansatz circuit. The dashed
box $U^{(k)}$ is the $k$th variational layer; $\theta^{(k)}_j$
($\phi^{(k)}_j$) is the $XX{+}YY$ ($ZZ$) angle shared by the $j$th
mirror-symmetric bond orbit in layer $k$ (odd $O_j\leftrightarrow
O_{M+1-j}$, even $E_j\leftrightarrow E_{M-j}$; $\theta=\phi$ at
$\delta=0$). Even-bond sublayer first. \label{fig:D06}}
\end{figure}

\begin{figure}[htbp]
\centering
\addtocounter{figure}{-1}
\includegraphics[width=\columnwidth]{exp04_D06b.pdf}
\caption{(b) Topological initial state plus required ansatz circuit, same
notation as (a); odd-bond sublayer first.}
\end{figure}

\begin{figure}[htbp]
\centering
\addtocounter{figure}{-1}
\includegraphics[width=\columnwidth]{exp04_D06c.pdf}
\caption{(c) AFM initial state plus required ansatz circuit, same notation
as (a); even-bond sublayer first.}
\end{figure}

\FloatBarrier
\section{Hardware implementation}
\label{sec:hardware}

\subsection{Premium-qubit selection}
\label{sec:premium}

To run on hardware we screen candidate chains (rings) with two
quantities: the stabilizer mean, probing two-qubit gate quality, and
the readout fidelity, probing measurement quality.
The ideal output is the cluster state
$|C\rangle=\prod_{\langle ij\rangle}CZ_{ij}|+\rangle^{\otimes n}$,
where $\langle ij\rangle$ runs over the edges of the candidate subgraph;
its stabilizers $S_i$ (two-body $X_0Z_1$, $Z_{n-2}X_{n-1}$ at open ends,
three-body $Z_{i-1}X_iZ_{i+1}$ otherwise) have expectation $+1$ in the
absence of errors, so deviations of $\langle S_i\rangle$ from $1$
directly reflect $CZ$ gate errors.
Each candidate is characterized by $4$ benchmark circuits: \texttt{stab\_g0}
and \texttt{stab\_g1} measure the even- and odd-index stabilizers separately
(the two groups anticommute and cannot be read out together), while
\texttt{allzero} and \texttt{allone} contain no $CZ$ gates and yield the readout
fidelity $F_{\rm ro}$ from the mean of $P(00\cdots0)$ and
$P(11\cdots11)$. The score
\begin{equation}
{\rm score}=0.8\,\bar{S}+0.2\,F_{\rm ro},
\end{equation}
where $\bar{S}=n^{-1}\sum_i\langle S_i\rangle$ is the stabilizer mean,
ranks the candidates, weighting two-qubit gate quality over readout.
Cluster states~\cite{raussendorf2001oneway}
are the ideal outputs whose stabilizers we measure; measurement-based
computation on such states is reviewed in
Ref.~\cite{briegel2009measurement}.
The $8$-chain and $10$-ring benchmarks differ (the ring carries one
extra closing $CZ$ and all-three-body stabilizers), so rankings are only
compared within each geometry, never across. In contrast to generic
randomized benchmarks, this screening probes exactly the stabilizers of
the states the hardware will prepare.
Hardware characterization around this screening takes several
complementary forms: Pauli-frame randomization has been demonstrated on
a superconducting qubit and validated by gate-set
tomography~\cite{ware2021experimental}, while the analysis of
randomized benchmarking under time-correlated dephasing shows where the
standard Markovian decay model stops
applying~\cite{qi2021randomized}. Shadow-based Hamiltonian measurement
offers a measurement-frugal
alternative~\cite{hadfield2022measurements,huang2021derandomization}, and
shadow tomography extends to quantum
channels~\cite{kunjummen2023shadow}.

The resulting rankings and selected qubits are reported in the Results and
Discussion feature. This Methods section defines only the screening
protocol and its within-geometry comparison rule.

\subsection{Premium-qubit rankings}
\label{sec:rankings}

The screening reports three rankings on the Baihua processor: (a) top-ten
$8$-chain candidates, (b) top-ten $10$-ring candidates, and (c) top-ten
$8$-subchains within the champion ring. The specific qubit tables follow
the latest S04 manifest; until the D07 figures land the slots below are
reserved and no ranking is asserted.

\begin{figure}[htbp]
\centering
\fbox{\parbox{0.9\columnwidth}{\centering TODO(GAP-001): D07a ranking figure pending.}}\\(a)
\fbox{\parbox{0.9\columnwidth}{\centering TODO(GAP-001): D07b ranking figure pending.}}\\(b)
\fbox{\parbox{0.9\columnwidth}{\centering TODO(GAP-001): D07c ranking figure pending.}}\\(c)
\caption{Premium-qubit ranking placeholders. \label{fig:D07}}
\end{figure}

\FloatBarrier
\subsection{Readout error mitigation}
\label{sec:readout}

Readout errors map the ideal distribution to the measured one,
\begin{equation}
\vec{P}_{\rm meas}=M\vec{P}_{\rm ideal},
\end{equation}
where $M$ is the assignment matrix with $M_{ab}=P(a|b)$.
Its columns are calibrated by preparation: preparing $|0\rangle$
gives the first column, preparing $|1\rangle$ the second,
\begin{equation}
M=\begin{bmatrix} P(0|0) & P(0|1) \\ P(1|0) & P(1|1) \end{bmatrix},
\end{equation}
where the diagonals read correct and the off-diagonals misread.
Mitigation inverts the map,
\begin{equation}
\vec{P}_{\rm mit}=M^{-1}\vec{P}_{\rm meas},
\end{equation}
where $\vec{P}_{\rm mit}$ is then projected onto the simplex by
truncating negatives to zero and renormalizing. In other words, the
raw counts are unfolded with the calibrated confusion matrix before any
observable is evaluated, following the tensor-product readout mitigation of
Bravyi~\textit{et al.}~\cite{bravyi2021mitigating}.
Assuming uncorrelated per-qubit errors, each qubit is corrected
independently at a calibration cost of $2$ circuits per qubit, linear in
system size; on hardware we recalibrate per batch against drift, with the
platform's built-in readout correction switched off.
The uncorrelated assumption keeps the overhead linear but cannot capture
crosstalk; correlated readout errors, if present, survive mitigation.
This tensor-product inversion sits within a broader mitigation
landscape. Probabilistic error cancellation expands noisy gates in a
quasiprobability basis and resamples
accordingly~\cite{temme2017error}, now extended to dynamic circuits with
mid-circuit measurement~\cite{gupta2024probabilistic} and unified with
noise scaling~\cite{mari2021extending}; zero-noise extrapolation instead
fits $E(\lambda)$ at boosted noise strengths and continues to
$E(0)$---first demonstrated on superconducting
hardware in Ref.~\cite{kandala2019error}---with optimized Richardson
protocols~\cite{krebsbach2022optimization}, inverted-circuit noise
estimation~\cite{koenig2024inverted}, and active error suppression by
artificial noise boosting inside variational
loops~\cite{li2017efficient}---though time-correlated noise is
systematically harder to extrapolate
away~\cite{schultz2022analyzing}. Matrix-free iterative unfolding scales
the same assignment-matrix idea to dozens of qubits without ever forming
$M$~\cite{nation2021scalable}; our per-qubit inversion is its minimal,
drift-robust endpoint.

\FloatBarrier
\section{Hardware-simulator comparison}
\label{sec:comparison}

We compare hardware and simulator on the same panels: simulator
$p=1,2,3$ curves overlaid with the hardware best-layer $p^*$ scatter
points, each estimator carrying its standard-deviation error bar.
That is, every point shows the sampling uncertainty of its own
estimator, so scatter across layers can be judged against shot noise
rather than read as physics.
Hardware and simulator are each internally normalized to $0$--$1$;
we compare trends, not absolute values. The separate normalization is
deliberate: gate-error suppression and readout residuals shift the two
platforms' absolute scales differently, while the $s$-dependence of the
AFM structure factor and the string operator survives as the common
signal.
Hardware variational experiments provide the backdrop: complete molecular
spectra from VQE with subspace expansion on two transmons~\cite{colless2018robust},
quantum chemistry on trapped
ions~\cite{hempel2018quantum}, hardware-efficient optimization of
six-qubit molecular and magnetic
problems~\cite{kandala2017hardware}, Hartree--Fock on twelve
superconducting qubits~\cite{google2020hartree}, and spin-chain ground
states up to $102$ qubits on cloud
devices~\cite{yu2023simulating}. Our protocol differs in comparing
hardware against simulator layer-by-layer on directly measurable
topological markers rather than energies alone.
Four panels cover $\delta=0$ and $\delta=0.85$ times the two
directly measurable markers; the figure slots are reserved below and the
trend statements stay frozen until the D08 figures land.

\begin{figure*}[htbp]
\centering
\begin{minipage}[b]{0.48\textwidth}
\centering
\fbox{\parbox{0.9\textwidth}{\centering TODO(GAP-002): D08a figure pending ($\delta=0$ AFM).}}\\(a)
\end{minipage}
\hfill
\begin{minipage}[b]{0.48\textwidth}
\centering
\fbox{\parbox{0.9\textwidth}{\centering TODO(GAP-002): D08b figure pending ($\delta=0$ string).}}\\(b)
\end{minipage}
\\
\begin{minipage}[b]{0.48\textwidth}
\centering
\fbox{\parbox{0.9\textwidth}{\centering TODO(GAP-002): D08c figure pending ($\delta=0.85$ AFM).}}\\(c)
\end{minipage}
\hfill
\begin{minipage}[b]{0.48\textwidth}
\centering
\fbox{\parbox{0.9\textwidth}{\centering TODO(GAP-002): D08d figure pending ($\delta=0.85$ string).}}\\(d)
\end{minipage}
\caption{Hardware-simulator comparison, four panels (placeholders):
(a) $\delta=0$ AFM structure factor, (b) $\delta=0$ string operator,
(c) $\delta=0.85$ AFM structure factor, (d) $\delta=0.85$ string
operator. Each panel will overlay simulator $p=1,2,3$ curves with
hardware $p^*$ scatter; normalization as in text. \label{fig:D08}}
\end{figure*}

\FloatBarrier
\section{Conclusion}
\label{sec:conclusion}

We introduced the dimerized XXZ chain and mapped its finite-size phase
diagram at $L=8$ with the normalized reflection marker
$\tilde{Z}_{\mathcal R}$, resolving trivial, topological, and AFM phases.
Gap scaling confirms finite-size effects shift only phase boundaries.
In contrast to markers that require full tomography, our hardware
protocol rests entirely on directly measurable correlators: variational
circuits with symmetry-preserving orbit ansatzes prepare the three
phases from matched initial states, premium-qubit screening selects the
best hardware subgraph, and readout mitigation unfolds the raw counts.
The protocol is deliberately modular---improved ansatzes, richer
screening metrics, or correlated readout models can each be upgraded
without changing the comparison framework---and it provides a modular
starting point for tracking topological markers on larger lattices.
The present study is limited in three ways: all phase assignments rest
on $L=8$ with gap scaling as the only size extrapolation; readout
mitigation assumes uncorrelated per-qubit errors; and the
hardware-simulator trend comparison awaits the D08 figures, until which
the hardware half of the comparison remains a reserved protocol rather
than a result.

\FloatBarrier
\appendix

\section{Finite-size gap analysis}
\label{sec:gap}

We take the gap in the $P=-2$ symmetry sector,
\begin{equation}
\Delta_{\rm sec}(L,s)=E_1^{(P=-2)}(L,s)-E_0^{(P=-2)}(L,s),
\end{equation}
where $E_0^{(P=-2)}$ ($E_1^{(P=-2)}$) is the ground-state (first-excited-state)
energy in that sector.
The sector gap at $\delta=0$ for $L=8,12,16$ is shown in
Fig.~\ref{fig:D02a}. Near criticality finite-size scaling gives
$\Delta_{\rm sec}L\approx A$; plotting $A$ versus $s$
(Fig.~\ref{fig:D02b}) the crossings for increasing $L$ stabilize at
$s=0.5$. Specifically, the drift of the crossing between successive sizes
shrinks as $L$ grows, consistent with convergence to $s=0.5$ in the
thermodynamic limit. The gap at $s=0.5$ versus $1/L$
($L=8,12,16,20,24\ $) is shown in Fig.~\ref{fig:D02c}; the linear fit through $L=20,24$ (coefficients
marked in the figure) has a small fitted intercept $b=0.05$, consistent with
the gap-closing trend. Finite-size effects therefore shift only the
phase boundaries slightly and leave the bulk phase structure
unaffected.
Ground states throughout are obtained by exact diagonalization
(Lanczos), and the crossing analysis follows standard finite-size
practice for quantum spin systems, for which we follow the lecture
notes of Ref.~\cite{sandvik2010computational}.

\begin{figure}[htbp]
\centering
\includegraphics[width=\columnwidth]{exp03_D02a.pdf}
\caption{Symmetry-sector gap $\Delta_{\rm sec}$ at $\delta=0$ for
$L=8,12,16$. \label{fig:D02a}}
\end{figure}

\begin{figure}[htbp]
\centering
\includegraphics[width=\columnwidth]{exp03_D02b.pdf}
\caption{$A=\Delta_{\rm sec}L$ versus $s$; crossings stabilize at
$s=0.5$ as $L\to\infty$. \label{fig:D02b}}
\end{figure}

\begin{figure}[htbp]
\centering
\includegraphics[width=\columnwidth]{exp03_D02c.pdf}
\caption{Gap at $s=0.5$ versus $1/L$ with linear fit through
$L=20,24$ (small fitted intercept $b=0.05$). \label{fig:D02c}}
\end{figure}

% CIT-003 (optional): finite-size scaling direction (bind)


\FloatBarrier
\bibliography{references}

\end{document}
